\section{统一之后的家庭、法律与区域断裂}
\label{sec:western-jin-households-law-regions}

二六五年司马炎受禅，二八〇年王濬舰队抵建业，西晋在形式上接收魏、蜀、吴三套州郡、官民与道路网络。统一并没有使这些网络自动使用同一套赋役、司法和地方关系。益州水军、荆州军府、淮南城戍与江东州县在伐吴中能够被组织为一项行动，战后如何重置官员、仓储、户籍和降附者，才决定“统一”在何处有具体内容。史书更擅长书写受禅、舰队和降仪，较少留下江南、巴蜀与北方家庭在接收中的连续记录。

\subsection{律令与地方秩序}

二六八年《泰始律》完成，试图整理魏晋以来刑律和审断框架。律令的颁行可以核对，却不能证明地方狱讼已同步使用统一程序。县廷、州府、宗族、部曲和有军事力量的地方人士都会影响案件如何被提出、调解或执行。法律文本首先说明中枢希望何种分类、责任和程序能够被书写；它不能自行显示贫富、籍属、性别或地域如何改变当事人的实际处境。把泰始律仅当作“法治进步”或“制度复古”，都会跳过它从中央条文抵达地方审断的复杂距离。

羊祜二六九年在襄阳经营屯戍、招纳和对吴关系，二七二年西陵督步阐降晋、陆抗收复峡口，说明同一法律与官制之外，边将、城戍和地方联络仍能决定江汉秩序。招纳的数字、将领品格和怀柔效果在《晋书》《三国志》之间各有修辞；可以确定的是，襄阳、峡口和武昌以西的道路、仓储和渡口必须被不断维护。晋军二七九年多路南下的胜利，依赖的正是此前这些并不均质的边防准备。

\subsection{继承危机为何会进入地方}

武帝死后，皇后、宗王、禁军和近侍围绕继承与辅政争夺中枢。八王之乱通常以宗室人物的联盟、背叛和死亡叙述，然而每一次调兵都需经过州郡的粮饷、道路和人事网络。宗王名义上的封国、都督权和实际能够调用的军队并不一致；地方官与军人也并非只按血缘或诏令选择归属。宫廷冲突因而会沿着黄河、关中、河北和江淮的资源节点外溢，反过来又使地方军队进入继承政治。

二九六年至二九九年齐万年在关中起兵，氐羌集团、关中军政和西晋处置相互卷入。史书族名不能替代部众内部差异，也不能将关中冲突写成某一族群天生反晋。战争暴露的是军费、饥馑、地方控制和士卒关系的脆弱，而非单一“民族矛盾”。关中一旦需要长期征发，原有的土地、迁徙和军事安排便会被重新打乱；这种压力后来同八王间的调兵交织，削弱了中枢再次稳定地区的能力。

\subsection{文书、墓葬与地方记忆}

二八一年汲郡魏襄王墓出竹书，朝廷获得一批来自墓葬的古代文本。竹书的发现、整理和讨论可由相关记述追踪，原书写年代、墓葬归属、整理损失以及后世利用必须分开处理。它提示西晋不只是战争与官制的时代：文字、墓葬、学术权威和朝廷政治也在相互作用。墓中所存文本并不能代表一般读写经验，却使人看见地方埋藏、士人解释与国家知识如何偶然接合。

《晋书》帝纪、志与列传、《三国志》陆抗等传及《资治通鉴》构成统一、律令和边防的年序；汲冢竹书传本、墓志、城址与区域考古提供另一种局部尺度。后出的《晋书》已站在统一破裂之后，人物品评与责任归属尤其需要警惕。本节因此不把西晋理解为从二八〇年胜利直线走向崩解的王朝，而是把法律、接收、边防和地方社会看作始终未被完全接合的条件。
